\section{Conclusions and future work}\label{sec:conclusions}

This paper introduced \textbf{Structural Cognition} as a new paradigm for steel structures and articulated its concrete embodiment, the \textbf{Self-Aware Steel Structure}: a structure able to perceive, process, interpret and continuously report its own mechanical state. The central argument is that conventional structural health monitoring, even when it delivers real-time measurement, largely stops at data acquisition and storage, leaving interpretation to periodic human analysis. Grounded in the statistical pattern recognition paradigm and the fundamental axioms of SHM \citep{farrar2013shm}, the proposed framework reframes the structure itself as the locus of interpretation. Rather than merely collecting strain, acceleration and temperature streams, a Self-Aware Steel Structure recognizes healthy behavior, identifies anomalies, detects stiffness loss and bolted-connection loosening, classifies damage severity, and autonomously communicates its condition along the Rytter hierarchy of existence, location, type, extent and prognosis.

The contribution is threefold. The \textbf{conceptual paradigm} distinguishes data-collecting SHM from a structure that interprets its own responses and maintains a permanent, bidirectional link to its digital representation. The \textbf{reference architecture} then layers a distributed multi-sensor network (strain gauges, triaxial accelerometers, temperature, displacement and inclinometer channels), a finite element baseline of static, modal, harmonic and transient behavior, a continuously updated digital twin, a machine-learning cognition layer, and a governed decision-and-reporting layer that issues autonomous assessments under explicit governance. This architecture is theoretically coherent: modal parameters depend only on mass and stiffness, so that damage can produce detectable shifts in natural frequencies, mode shapes and damping when excitation and observability are adequate \citep{chopra2017dynamics,cook2002fea}; data-driven dimensionality reduction and model discovery compress distributed signals into low-dimensional damage-sensitive features and a healthy baseline \citep{brunton2019data}; and supervised classifiers (Random Forest, SVM, ANN), unsupervised novelty detectors (autoencoders) and self-supervised temporal models (LSTM) map onto the SHM division between detecting damage and characterizing its type, severity and evolution \citep{farrar2013shm,naser2023ml,wangcha2022unsupervised}. Finally, the \textbf{six-step methodology} sequences FEM development, physical construction, sensor installation and calibration, static and dynamic testing, digital-twin integration, and AI training, validation and deployment into a structured, auditable engineering protocol.

The expected impact spans the lifecycle of steel infrastructure. By converting passive load-resisting members into self-interpreting systems, the framework promises improved \textbf{safety} through continuous, near-real-time condition screening; reduced \textbf{inspection cost} by shifting from time-based to condition-based and predictive maintenance; and increased \textbf{reliability} through traceable, physics-consistent damage indicators. Together these elements point toward \textbf{cognitive digital twins} and autonomous predictive-maintenance systems for the next generation of intelligent structures \citep{chiachio2022structural,zou2026deep}.

We emphasize that the methodology, testbed and AI layer described here are \textbf{planned}; no experiments have yet been performed, and the framework is offered as a conceptual and architectural proposal rather than a validated result. The concrete next steps are the following:

\begin{enumerate}
    \item \textbf{Build the three-dimensional braced SHS frame demonstrator} with bolted, semi-rigid connections, and instrument it at sensor positions selected by preliminary numerical analysis.
    \item \textbf{Develop and calibrate the finite element model}, establishing baseline modal and harmonic behavior and verifying it against independent solutions and experiment \citep{cook2002fea}.
    \item \textbf{Run the planned static and dynamic tests} using the variable-speed rotating-unbalance excitation to traverse natural frequencies and induce controlled, reversible stiffness changes.
    \item \textbf{Integrate the digital twin} with live sensor feeds and quantify deviations via statistical indicators.
    \item \textbf{Train, validate and deploy the AI cognition layer}, using rigorous metrics and cross-validation, addressing scarce damaged-state data through FE-augmented datasets, and incorporating explainability so that autonomous integrity assessments remain physically interpretable and trustworthy \citep{naser2023ml}.
\end{enumerate}

Demonstrating closed-loop structural cognition on this testbed is the immediate objective of future work.
